% theme: quadratic_proportional_function_meaning_and_formula
\MajorQuestionLesson{二乗に比例する関数の意味と式}
\SetRowSpace{6mm}

\TypeQuestionDouble{次の数量関係について、$y$ が $x$ の二乗に比例する関数である場合は○、そうでない場合は×を答えなさい。}{hantei}
\QQRow{$y=4x^2$}{○}{$y=4x$}{×}
\QQRow{$y=-3x^2$}{○}{$y=x^2+5$}{×}
\QQRow{$y=\dfrac{1}{2}x^2$}{○}{$y=\dfrac{12}{x}$}{×}

\TypeQuestionSingle{次の数量関係について、$y$ を $x$ の式で表しなさい。}{bunsho}
\QQ{一辺が $x\mathrm{cm}$ の正方形の面積を $y\mathrm{cm}^2$ とする。}{$y=x^2$}
\QQ{半径が $x\mathrm{cm}$ の円の面積を $y\mathrm{cm}^2$ とする。}{$y=\pi x^2$}
\QQ{底辺と高さがともに $x\mathrm{cm}$ の三角形の面積を $y\mathrm{cm}^2$ とする。}{$y=\dfrac{1}{2}x^2$}
\QQ{正方形の土地の一辺を $x\mathrm{m}$、土地代を $y$ 円とする。$1\mathrm{m}^2$ あたりの土地代は $8000$ 円である。}{$y=8000x^2$}

\TypeQuestionSingle{次の式について、比例定数を答えなさい。}{hireiteisu}
\QQRow{$y=5x^2$}{$5$}{$y=-2x^2$}{$-2$}
\QQRow{$y=\dfrac{3}{4}x^2$}{$\dfrac{3}{4}$}{$y=-\dfrac{1}{6}x^2$}{$-\dfrac{1}{6}$}
\QQRow{$y=x^2$}{$1$}{$y=-x^2$}{$-1$}

\LessonPageBreak

\TypeQuestionSingle{$y$ が $x$ の二乗に比例し、次の条件を満たす。比例定数 $a$ を求めなさい。}{akettei}
\QQRow{$x=2$ のとき $y=12$}{$a=3$}{$x=3$ のとき $y=-18$}{$a=-2$}
\QQRow{$x=-4$ のとき $y=8$}{$a=\dfrac{1}{2}$}{$x=-5$ のとき $y=-75$}{$a=-3$}
\QQRow{$x=6$ のとき $y=9$}{$a=\dfrac{1}{4}$}{$x=-2$ のとき $y=-1$}{$a=-\dfrac{1}{4}$}

\TypeQuestionSingle{$y$ が $x$ の二乗に比例し、次の条件を満たす。$y$ を $x$ の式で表しなさい。}{shikikettei}
\QQRow{$x=2$ のとき $y=20$}{$y=5x^2$}{$x=3$ のとき $y=6$}{$y=\dfrac{2}{3}x^2$}
\QQRow{$x=-2$ のとき $y=-12$}{$y=-3x^2$}{$x=-4$ のとき $y=4$}{$y=\dfrac{1}{4}x^2$}
\QQRow{$x=5$ のとき $y=-10$}{$y=-\dfrac{2}{5}x^2$}{$x=-6$ のとき $y=12$}{$y=\dfrac{1}{3}x^2$}

\TypeQuestionDouble{$y$ が $x$ の二乗に比例し、次の条件を満たす。指定された $x$ の値に対応する $y$ の値を求めなさい。}{gyakusan}
\QQ{$x=2$ のとき $y=8$ である。$x=5$ のときの $y$}{$50$}
\QQ{$x=3$ のとき $y=-27$ である。$x=-2$ のときの $y$}{$-12$}
\QQ{$x=-4$ のとき $y=12$ である。$x=8$ のときの $y$}{$48$}
\QQ{$x=-6$ のとき $y=-18$ である。$x=2$ のときの $y$}{$-2$}
